\subsection{Every one-step short-one terminal with long arm at least four}\label{proof:Q5-C038}
For every finite $b\ge4$, keep the actual orientation $F_b=qf_b$,
one connector toward the receiving root, and its outward length-two
arm $f_2$. We prove
\begin{equation}\tag{L38.0.1}\label{eq:L38:claim}
\E[f_2\{T(qf_b)\}^4]\le1-(47/65000)d.
\end{equation}
All quantities use the original finite measure. Write $c=\kappa=19/32$,
$r=Kg$, and $\delta_j=\|K^j\one\|_2^2$.
The original row constraint gives $|r|\le pq/2$ and
$0\le f_2\le2-q^2/2$.
Put
$R=\sum_{j=2}^bK^j\one$, $\beta=\beta_b$,
$a=1-\beta/2$, $t_0=\beta_{b+1}=-\E R$,
$k_0=\E KR=d-\beta_{b+2}$, $V=\E R^2$.
Actual positive convolution retains its $\beta_bf_0$ term and gives
$q+f_b\le2a+R$. Thus
\begin{equation}\tag{L38.2.1}\label{eq:L38:envelope}
0<T(qf_b)\le Z:=a^2q-at_0-aKR+V/2.
\end{equation}
For $d>0$, $v=\delta_2/d\in[d,\theta]$ and
$m:=1-d/2\le a\le A_0:=1-(1-v)d/2$, $5/6\le a\le1$.

For every finite $N\ge3$ the exact geometric polynomial satisfies
\begin{equation}\tag{L38.2.3}\label{eq:L38:finiteG}
|G_N(i\sqrt x)|^2\le1-(7/15)x\quad(0\le x\le\theta).
\end{equation}
At $x=0$ equality holds. For odd $N$, the finite formula reduces to
$x^{N-1}+(7/15)x\le8/15$, whose left side is at most
$\theta^2+(7/15)\theta=5447/15360<8/15$.
For $N\equiv0\pmod4$, the numerator $(1-x^{N/2})^2\le1$ and
$(7/15)(1+x)\le21/32<1$.
For $N\equiv2\pmod4$, $N\ge6$, it suffices that
$2x^{N/2-1}+x^{N-1}+(7/15)x\le8/15$.
Use $8/15-2\theta^2-(7/15)\theta=7/512$ and
$\theta^5=371293/33554432<7/512$ (numerators $371293<458752$).
This covers every finite integer.
Consequently, since $b-1\ge3$,
\[
\begin{gathered}
V\le\delta_2-(7/15)\delta_3,\quad
U:=\E(KR)^2\le\delta_3-(7/15)\delta_4\le\theta_*\delta_2,\\
\theta_*=5057/15360<1/3,\quad\theta-\theta_*=1183/15360>1/13.
\end{gathered}
\]
The last endpoint bound follows from increasing $x-7x^2/15$ on
$[0,\theta]$, and $1183\cdot13-15360=19$.
Set $\rho=\E(rR)$ and $\sigma=\E(sR)$.
The exact finite signed sums are
\[
\begin{aligned}
\rho&=\sum_{h=0}^{\lfloor b/2\rfloor-1}(-1)^h\delta_{h+2},\\
k_0&=\sum_{h=0}^{\lfloor(b-1)/2\rfloor-1}(-1)^h\delta_{h+2},\quad
\sigma=\sum_{h=0}^{\lfloor(b-1)/2\rfloor-1}(-1)^h\delta_{h+3}.
\end{aligned}
\]
Thus $c\delta_2\le\rho,k_0\le\delta_2$.
For $b=2h\ge4$, $\rho+\sigma=\delta_2$, $t_0+k_0=d$;
for $b=2h+1\ge5$,
$\rho+\sigma=\delta_2+(-1)^{h-1}\delta_{h+2}$ and
$t_0+k_0=d+(-1)^{h+1}\delta_{h+1}$.
The first positive extra term is at $b=7$, so with
$L_0=1+\theta^3=34965/32768$,
\begin{equation}\tag{L38.2.8}\label{eq:L38:jointmeans}
\rho+\sigma\le L_0\delta_2,\qquad t_0+k_0\ge d-\delta_3.
\end{equation}

The normalized positive bound $\widehat Z=Z/a^2$ satisfies
$\widehat Z\le q+(4/5)\sqrt d-3d/8$,
$\E\widehat Z\le1-5d/8$.
Indeed $k=1/a\in[1,6/5]$, $V\le\delta_2$,
$k\sqrt\theta\le4/5$, and
$kt_0-k^2V/2\ge(ck-\theta k^2/2)d$.
The last concave quadratic has endpoint values $25/64,21/50$,
both at least $3/8$.
The normalized mean loss is at least
$[k(1-\theta^2)-k^2\theta/2]d$, whose endpoint values
$647/1024,9081/12800$ exceed $5/8$.
For any nonnegative $z,u_0$ satisfying these same bounds, the full
continuous coefficient is
\begin{equation}\tag{L38.3.3}\label{eq:L38:curvature}
f_2(z^2+2u_0z+3u_0^2)\le29/3.
\end{equation}
To prove it put $x=\sqrt d$, $n=1-5x^2/8$,
$t=1+4x/5-x^2$. The bracket is at most
$(q+t)^2+2n^2$, whose exact expansion is
\[
q^2+2q+3+(8/5)(q+1)x-(2q+193/50)x^2
-(8/5)x^3+(57/32)x^4.
\]
Since $x<3/5$, its cubic and quartic sum is nonpositive.
Complete the square in the remaining quadratic to add at most
$32(q+1)^2/(100q+193)$ to $q^2+2q+3$.
The base after multiplication by $2-q^2/2$ is at most nine.
The added part is below $2/3$ because, for $y=q-1\in[-1,1]$,
\[
100q+193-24(4-q^2)(q+1)^2
=5+4y+24y^2(y+3)^2\ge1.
\]
All denominators are positive, including both $q$ endpoints.
By homogeneity the actual $Z$, $\mu_0=\E Z$ has coefficient
$Ca^4$, $C=29/3$.

For $w=Z-\mu_0$, exact original skew calculations give
\begin{equation}\tag{L38.4.1}\label{eq:L38:exact}
\begin{gathered}
\mu_0=a^2-a(t_0+k_0)+V/2,\quad \E[f_2w]=-a^2d+a(\rho+\sigma)-adk_0,\\
\E w^2=a^4d-2a^3\rho+a^2(U-k_0^2).
\end{gathered}
\end{equation}
The actual mean square $k_0^2$ is retained.
Equations \eqref{eq:L38:finiteG}--\eqref{eq:L38:jointmeans} imply
\begin{equation}\tag{L38.4.3}\label{eq:L38:M}
\mu_0\le M(d,v):=1-[2-(1739/960)v]d
+[(1-v)(3-29v/16)/4]d^2.
\end{equation}
Before taking $a\le A_0$ the bound is
$a^2-ad+(a-7/30)\delta_3+\delta_2/2$;
$a-7/30\ge3/5$ permits $\delta_3\le\theta\delta_2$,
and its derivative in $a$ is $2a-d+\theta\delta_2>0$.
The weighted centered mean is at most $dH(d,v)$, where
\[
H=-m^2+m(L_0-cd)v<0.
\]
The upper expression decreases in $a$, derivative
$-2a+(L_0-cd)v<0$, justifying $a\ge m$.
The original variance and homogeneous curvature give
\begin{equation}\tag{L38.4.5}\label{eq:L38:B}
Ca^4\E w^2\le B_*(d,v):=CdA_0^6Q_*,\quad
Q_*=A_0^2-(2A_0c-\theta_*)v-c^2dv^2.
\end{equation}
The bracket increases in $a$ and upper-bounds a nonnegative actual
variance; on the full domain it is also positive, since
$A_0\ge5/6$, $2c\theta<1/2$, $c^2dv^2<1/48$ yield
$Q_*>25/144$.
The exact fourth expansion therefore gives the two-endpoint bound
\begin{equation}\tag{L38.4.7}\label{eq:L38:endpoints}
\E[f_2Z^4]\le\max\{B_*,(1-d)N^4+4dHN^3+B_*\}
\quad(\mu_0\le N).
\end{equation}
This follows because $t\mapsto(1-d)t^4+4dHt^3$ decreases then
increases, or increases throughout, on $[0,N]$.

The actual envelope has \(\mu_0=\E Z\ge0\), by
\eqref{eq:L38:envelope}. Take \(N=M(d,v)\) from \eqref{eq:L38:M}.
On the whole moment triangle \(0\le d\le1/3\), \(d\le v\le\theta\),
this bound satisfies \(M\ge1-2d\ge1/3\): its remaining displayed
terms are nonnegative. Moreover
\(H\le m[-5/6+L_0\theta]<0\), where \(m=1-d/2\).
For \(\phi(t)=(1-d)t^4+4dHt^3\),
\(\phi'(t)=4t^2[(1-d)t+3dH]\).
The bracket increases, so \(\phi\) decreases and then increases, or
is monotone, on \([0,M]\). Its maximum is at zero or \(M\).
Thus \eqref{eq:L38:endpoints} reduces the actual integral to
\[
 \max\{B_*,\ (1-d)M^4+4dHM^3+B_*\}.
\]
Put \(\lambda=47/65000\), and define the two polynomial gaps
\begin{equation}\label{eq:reference-synthesis:C038-gaps}
\begin{aligned}
 \Phi(d,v)&=(1-d)M(d,v)^4+4dH(d,v)M(d,v)^3+B_*(d,v),\\
 D_0(d,v)&=1-\lambda d-B_*(d,v),\\
 D_1(d,v)&=[1-\lambda d-\Phi(d,v)]/d.
\end{aligned}
\end{equation}
The second numerator vanishes identically at \(d=0\); after exact
cancellation, \(D_1\) is a polynomial there as well.
The substitutions
\[
 d=x/3,\qquad v=d+(\theta-d)y,\qquad (x,y)\in[0,1]^2
\]
cover the full triangle, since \(\theta-d\ge7/96>0\).
They may include moment pairs that no kernel realizes; proving the
polynomial bounds on this larger set is sufficient.
Put \(\widetilde D_i(x,y)=D_i(x/3,x/3+(\theta-x/3)y)\).
These two polynomials have the following complete-matrix Bernstein bounds:
\begin{table}[htbp]\centering\small
\caption{Joint-parameter payments for the genuine one-step reference.}
\label{tab:reference-synthesis:C038}
\begin{tabular}{@{}cccc@{}}
\toprule
Polynomial&bidegree&number of coefficients&coefficient lower bound\\
\midrule
\(\widetilde D_0\)&\((17,8)\)&162&\(1/8\)\\
\(\widetilde D_1\)&\((16,8)\)&153&\(1/2\)\\
\bottomrule
\end{tabular}
\end{table}
Each entry is obtained from the displayed \(M,H,B_*\) by
\eqref{eq:bernstein-coefficients}.
All clearing denominators are positive constant integers.
Lemma~\ref{lem:bernstein-certificate} gives \(\widetilde D_0\ge1/8\)
and \(\widetilde D_1\ge1/2\) on the full closed square, hence the
same gap bounds throughout the original triangle.
For \(d>0\) these inequalities make both actual endpoint budgets
strictly smaller than \(1-\lambda d\), proving \eqref{eq:L38:claim}.
At $d=0$ original leaf integration gives equality without $K=0$.
Actual equal-pair decorations decrease the positive messages; full
reversal simultaneously reverses every original edge and its weight.
